\documentclass[10pt,a4paper]{amsart}
\usepackage[margin=19mm]{geometry}
\usepackage{amsmath,amssymb,mathtools}
\usepackage[scr=rsfs,cal=euler]{mathalfa}
\usepackage{xfrac}
\usepackage[unicode,hidelinks,bookmarks=false]{hyperref}
\newcommand{\ie}{i.e.\ }
\newcommand{\cf}{cf.\ }
\newcommand{\resp}{respectively\ }
\newcommand{\Subsection}{\subsection}
\providecommand{\nobreakdash}{\nobreak\hbox{-}}
\setlength{\emergencystretch}{2em}
\multlinegap0pt
\usepackage{amssymb}
\usepackage{float}
\usepackage[normalem]{ulem}
\usepackage{xcolor}
\newtheorem{theorem}{Theorem}
\newcommand{\Fq}{\mathbb{F}_{q}}
\DeclareMathOperator{\lcm}{lcm}
\title{\LARGE \bfseries Equivalence of Families of Polycyclic Codes over Finite Fields}
\author{Hassan Ou-azzou, Anna-Lena Horlemann}
\date{Source: arXiv:2503.04498v3 (CC BY 4.0)}
\begin{document}
\maketitle

\noindent\emph{Source and licence.} \LARGE \bfseries Equivalence of Families of Polycyclic Codes over Finite Fields. Version arXiv:2503.04498v3 (CC BY 4.0), distributed by arXiv under the Creative Commons Attribution 4.0 International licence, http://creativecommons.org/licenses/by/4.0/, as recorded in the arXiv licence metadata for this exact version and shown on the abstract page https://arxiv.org/abs/2503.04498v3. Full licence notice: \texttt{licenses/arxiv-2503.04498-CC-BY-4.0.txt}.\par\smallskip

\noindent\textit{Editorial scope.} Counted unit: the article's theorem Th\_general (source0.tex lines 1029--1045). The article's complete original proof of it is delivered with it (source0.tex lines 1047--1104, one \texttt{proof} environment of 58 lines). No mathematics is written, repaired or re-derived: the statement and the proof are the article's own lines, in the article's own notation, and no retained line is substituted or re-flowed.  Retained uncounted as the source-local context the proof needs: lines 376--395.  Nothing was omitted from any retained span, so the package carries no omission record.  No retained line contains a citation, so the wrapper carries no bibliography and no bibliography entry.  No retained line contains a figure environment, an image-inclusion command or drawing code, so no image is delivered and the package declares no asset dependency.  Editorial formatting only, in addition: an AMS \texttt{amsart} preamble that restates the article's own declarations and macros used by the retained lines, namely the environments \texttt{theorem} on separate counters, together with the standard packages those lines need.  The retained environments print in this document as Theorem 1 in order of appearance, whereas the article prints the same items with the numbering of its own full document.  The complete native source of the article as distributed in the arXiv e-print for this exact version is bundled unchanged as \texttt{sources/174334/source0.tex}.
\begin{theorem}\label{Th.1}
Let $0<\ell<n$ be an integer,  $(a_0, a_{\ell})$ and $(b_0, b_{\ell})$ be elements of $\mathbb{F}_q^{*} \times \mathbb{F}_q^{*}$,  and  $\xi$ be  a primitive element of $\Fq$. 
The following statements are equivalent:

\begin{enumerate}
    \item $(a_0, a_{\ell}) \sim_{(n,\ell)} (b_0, b_{\ell}).$

    \item The polynomials $ a_i x^{n-i } - b_i \in \mathbb{F}_q[x]$, with $i \in \{0, \ell\}$, have a common root in $\mathbb{F}_q^*.$

    \item The polynomial $\gcd(a_0 x^n - b_0, a_{\ell} x^{n - \ell} - b_{\ell})$ has at least one root in $\mathbb{F}_q^*.$
       \item The polynomial $\gcd(x^n - b_0 a_0^{-1}, x^{n - \ell} - b_{\ell} a_{\ell}^{-1})$ has at least one root in $\mathbb{F}_q^*.$
     \item  There exists $\alpha\in \Fq^*$ such that $ (a_0,a_{\ell})\star ( \alpha^n,   \alpha^{n-\ell})= (b_0,b_{\ell}) $.
     %,$ where $ \star$ is the component-wise product on $ \Fq^*\times \Fq^*. $ 
     
      \item  $ (a_0, a_{\ell})^{-1} \star(b_0,b_{\ell})\in H,$ where   $H$ is the cyclic subgroup of $\Fq^*\times \Fq^*$ generated by $(\xi^n,\xi^{n-\ell}) .$
\end{enumerate}
The equivalence between (1) and (6) implies that the number of $(n,\ell)$-equivalence classes is $$N_{(n,\ell)}:= \dfrac{ (q-1)^2}{ \lcm(\frac{q-1}{ \gcd(n,q-1)}, \frac{q-1}{ \gcd(n-\ell,q-1)})} = (q-1) \gcd(q-1,n,\ell).
%(q-1)\gcd\left( \frac{q-1}{\gcd(n,q-1)}, \frac{q-1}{\gcd(n-\ell ,q-1)} \right).
$$
\end{theorem}

\begin{theorem}\label{Th_general}
Let $ \vec{a}=(a_0,a_1,\ldots, a_{n-1}), \vec{b}=(b_0,b_1,\ldots, b_{n-1})\in \mathbb{F}_q^{n}$ have non-zero entries in the same $m$ positions, i.e., $a_{i_j}$ and $b_{i_j}$ are non-zero  for $0\leq i_0<\dots<i_{m-1}\leq n-1$. Moreover, let $\xi $ be a primitive element of $\Fq$. Then the following statements are equivalent:
\begin{enumerate}
\item $ \vec{a} \sim_{n} \vec{b}$. 

\item The polynomials $ a_{i_j} x^{n-i_j} - b_{i_j} \in \mathbb{F}_q[x]$, for $ j \in \{0,1,\ldots,m-1\}$, have a common root in $\mathbb{F}_q^*$.

\item The polynomial $ \gcd( a_{i_0} x^{n-i_0} - b_0, a_{i_1} x^{n-i_1} - b_{i_1}, \ldots, a_{i_{m-1}} x^{n-i_{m-1}} - b_{i_{m-1}})$ has at least one root in $\mathbb{F}_q^*$.

\item The polynomial $ \gcd_{\{0 \leq j \leq m-1\}}( x^{n-{i_j}} - b_{i_j} a_{i_j}^{-1}) $ has at least one root in $\mathbb{F}_q^*$.
 \item  There exists $\alpha\in \Fq^*$ such that $$ (a_{i_0}, a_{i_1},\ldots, a_{i_{m-1}})\star(\alpha^{n-{i_0}}, \alpha^{n-{i_1}} , \ldots, \alpha^{n-{i_{m-1}}})=  ( b_{i_0}, b_{i_1},\ldots, b_{i_{m-1}}).$$
\item $ (a_{i_0}, a_{i_1},\ldots, a_{i_{m-1}})^{-1}\star( b_{i_0}, b_{i_1},\ldots, b_{i_{m-1}}) \in H,$ where $H$ is the cyclic subgroup of $ (\Fq^*)^m$ generated by 
$( \xi^{n-{i_0}},\xi^{n-{i_1}},\ldots, \xi^{n-{i_{m-1}}} ).$
\end{enumerate}
In particular the number of $n$-equivalence classes is
$$ N=\dfrac{ (q-1)^m}{ \lcm_{ (0 \leq j\leq m-1)}\left(\frac{q-1}{ \gcd(n-i_j,q-1)}\right)}. $$ 
\end{theorem}

\begin{proof}
\begin{enumerate}
    \item [(1) $\Rightarrow$ (2)] 
    Suppose that $ \vec{a} \sim_{n} \vec{b} $, then there is $\alpha \in \mathbb{F}_q^*$ such that

$$\varphi_{\alpha} : \mathbb{F}_q[x] /\langle x^n - \vec{b}(x) \rangle \to \mathbb{F}_q[x] /\langle x^n - \vec{a}(x) \rangle, \quad f(x) \mapsto f(\alpha x) $$
is an $\mathbb{F}_q$-algebra isometry.  It follows that

$$ 
\varphi_{\alpha}(x^k) =  \alpha^k x^{k} \mod (x^n - \vec{a}(x)), \quad \forall  \ k = 0, 1, \ldots, n-1.
$$

As $\varphi_{\alpha}$ is an $\mathbb{F}_q$-algebra isometry and $ \varphi(x^n - \vec{b}(x)) = 0 \mod (x^n - \vec{a}(x))$, then
\begin{equation}
\varphi_{\alpha}(x^n) = \varphi_{\alpha}(\vec{b}(x)) = b_0 + \alpha b_1 x + \ldots + \alpha^{n-1}b_{n-1} x^{n-1}, \mod (x^n - \vec{a}(x)).
\end{equation}

On the other hand,
\begin{equation}
\varphi_{\alpha}(x^n) = \alpha^n x^n = \alpha^n ( a_0 + a_1 x + \ldots + a_{n-1} x^{n-1}), \mod (x^n - \vec{a}(x))
\end{equation}

Comparing term by term, we deduce that for any $ i \in \{0, 1, \ldots, n-1\}, \ a_i \alpha^{n-i} = b_i$, which means that $\alpha$ is a common root of the polynomials $ a_i x^{n-i} - b_i $, for $ i \in \{0, 1, \ldots, n-1\}$. As  $ a_{i_j}$'s and $  b_{i_j}$'s are the non-zeros components of $\vec{a}$ and $\vec{b},$   then $ a_{i_j} \alpha^{n-{i_j}} = b_{i_j},\ j=0,\ldots,m-1$.

\item[(2) $\Rightarrow$ (3)] and  (3) $\Rightarrow$ (4) are immediate. 

 \item [(4) $\Rightarrow$ (5)]
    Let $\alpha $ be a root of the polynomial   $\gcd_{\{0 \leq i \leq n-1, \ a_i \neq 0\}}( x^{n-i} - b_i a_i^{-1}) $. Then $\alpha$ is a common root of the polynomials $  x^{n-i_j} - b_{i_j} a_{i_j}^{-1} \  \text{for any} \ j\in \{ 0,1,\ldots, m-1\},$ and so   $a_{i_j} \alpha^{n-i_j}= b_{i_j}$.
    It follows that 

$$  (b_{i_0}, b_{i_1},\ldots, b_{i_{m-1}} )=  (\alpha^{n-i_0}, 
\alpha^{n-i_1} , \ldots, \alpha^{n-i_{m-1}})\star( a_{i_0},  a_{i_1}, \ldots, a_{i_{m-1}}).$$ 
\item [(5) $\Rightarrow$ (6)]
Suppose that there is $\alpha\in \Fq^{*}$ such that 
$$  (b_{i_0}, b_{i_1},\ldots, b_{i_{m-1}} )=  (\alpha^{n-i_0}, 
\alpha^{n-i_1} , \ldots, \alpha^{n-i_{m-1}})\star( a_{i_0},  a_{i_1}, \ldots, a_{i_{m-1}}).$$ 
For $\alpha= \xi^h,$ we obtain that 
$$  ( a_{i_0}, a_{i_1},\ldots, a_{i_{m-1}})^{-1}\star(b_{i_0}, b_{i_1},\ldots, b_{i_{m-1}} ) =   (\xi^{n-{i_0}}, \xi^{n-{i_1}} , \ldots, \xi^{n-i_{m-1}})^h .$$
It follows that $  ( a_{i_0},  a_{i_1}, \ldots, a_{i_{m-1}})^{-1}\star( b_{i_0},  b_{i_1}, \ldots, b_{i_{m-1}}) $ belongs to the cyclic subgroup  $H$ of $ (\Fq^{*})^m$ generated by $ (\xi^{n-{i_0}}, \xi^{n-{i_1}} , \ldots, \xi^{n-{i_{m-1}}}) .$
    \item [(6) $\Rightarrow$ (1)]  
     Suppose that $ (a^{n-{i_0}}, a^{n-{i_1}} , \ldots, a^{n-i_{m-1}})^{-1}\star(b^{n-{i_0}}, b^{n-{i_1}} , \ldots, b^{n-i_{m-1}}) $ is an element of the cyclic subgroup  $H$ of $ (\Fq^{*})^m$ generated by
$ (\xi^{n-{i_0}}, \xi^{n-{i_1}} , \ldots, \xi^{n-{i_{m-1}}}) $. Then there exists an integer $ h$ such that 
     $$  ( a_{i_0},  a_{i_1}, \ldots, a_{i_{m-1}})^{-1}\star( b_{i_0},  b_{i_1}, \ldots, b_{i_{m-1}}) =  (\xi^{n-{i_0}}, \xi^{n-{i_1}} , \ldots, \xi^{n-{i_{m-1}}})^h $$
    For $ \beta= \xi^{h}, $ we obtain  that $  a_j \beta^{n-j}= b_j, $ for any $ j\in \{ i_0,i_1,\ldots, i_{m-1}\}.$
As in the proof of Theorem \ref{Th.1}, we verify that  $ \varphi_{\beta}, $ as follows: 
  \begin{equation}
        \begin{array}{cccc}
        \varphi_{\beta} : & \mathbb{F}_q[x]/\langle x^n - \vec{b}(x) \rangle,  & \longrightarrow & \mathbb{F}_q[x] /\langle x^n - \vec{a}(x) \rangle, \\ 
        & f(x) & \longmapsto & f(\beta x),
        \end{array}
    \end{equation} 
	
	is an $\Fq$-algebra isometry with respect to the Hamming distance. 
\end{enumerate}
\noindent By the equivalence between (1) and (6), we deduce that the number of $n$-equivalence classes on on $(\Fq^*)^m$ corresponds to the order of the group $ (\Fq^*)^m/H,$ which equals $$ N=\dfrac{ (q-1)^m}{ \lcm_{ (0 \leq j\leq m-1)}\left(\frac{q-1}{ \gcd(n-i_j,q-1)}\right)}. $$ 

\qed
\end{proof}

\end{document}
